\documentclass[a4paper]{article}
\title{MA2: $(\sgn \sin n)^{\infty}_{n=1}$}
\author{Tommy Chu}
\date{}

\usepackage[czech]{babel}
\usepackage[utf8]{inputenc}
\usepackage[T1]{fontenc}
\usepackage{graphicx}
\usepackage{amsmath}
\usepackage{amsthm}
\usepackage{amssymb}
\usepackage{subfiles}
\usepackage{hyperref}
\usepackage{geometry}
\usepackage{mathtools}
\usepackage{algpseudocode}
\usepackage{algorithm}
\usepackage{physics}
\usepackage{dsfont}
\usepackage{bm}
\usepackage{bbm}
\usepackage{float}
\usepackage{enumitem}
\usepackage{graphicx}
\usepackage{multirow}
\usepackage{tikz}
\usepackage{tikz-cd}
\usepackage{pgfplots}
\usepackage{lmodern}
\usepackage{import}
\usepackage{microtype}
\usepackage{fancyhdr}
\usepackage{parskip}
\usepackage{paralist}
\usepackage{verbatim}
\usepackage{titling}

\graphicspath{ {./images/} }

\microtypesetup{
    protrusion=true,
    activate={true,nocompatibility},
    final,
    tracking=true,
    kerning=true,
    spacing=true,
    factor=1100,
}

\relpenalty   = 10000
\binoppenalty = 10000

\setlength{\droptitle}{-5em}

\newtheorem{theorem}{Tvrzení}
\DeclarePairedDelimiter\ceil{\lceil}{\rceil}
\DeclareMathOperator{\sgn}{sgn}

\begin{document}
\maketitle
\section*{Zadání}

Zjistěte, jestli je posloupnost $(\sgn \sin n)^{\infty}_{n=1}$ periodická.

Periodickou posloupností se myslí taková posloupnost $(a_n)^{\infty}_{n=1}$, pro které existuje perioda $T \in \mathbb{N}$ taková, že pro všechna $n \in \mathbb{N}$ platí $a_n = a_{n + T}$.

Funkce $\sgn : \mathbb{R} \rightarrow \{-1, 0, 1\}$ zobrazuje záporná čísla na -1, nulu na 0 a kladná čísla na 1.

\subsection*{Pomocná tvrzení}

\begin{theorem}
    Pokud má posloupnost $(\sgn \sin n)^{\infty}_{n=1}$ periodu $T \in \mathbb{N}$, pak ${T \geq 7 > 2\pi}$.
\end{theorem}
\begin{proof}
    $T \geq 6$  plyne přímo z výčtu prvních prvků: \[1, 1, 1, (-1), (-1), (-1), 1, \cdots\]
    $T = 6$ lze vyvrátit protipříkladem: ${\sgn \sin(1 + 5 \cdot T) = -1 \neq 1 = \sgn \sin(1)}$.
\end{proof}

\begin{theorem}
    Nechť je $a \in \mathbb{N}$. Neexistuje $k \in \mathbb{Z}$, pro které platí $a = k\pi$.
\end{theorem}
\begin{proof}
    Platí triviálně pro $k = 0$. Pro $k \neq 0$ by muselo platit $\pi = \frac{a}{k} \in \mathbb{Q}$, což je ve sporu s tím, že $\pi$ je iracionální číslo.
\end{proof}

\subsection*{Řešení}

Posloupnost $(\sgn \sin n)^{\infty}_{n=1}$ není periodická.

\begin{proof}
    Pro spor předpokládejme, že posloupost $(a_n)^{\infty}_{n=1}$, kde $a_n = \sgn \sin n$, je periodická s~periodou $T \in \mathbb{N}$.

    Zřejmě platí $\sin 1 > 0$, tedy $a_1 = 1$. Podle předpokladu dále platí $a_1 = a_{1 + T}$. Pro $a_{1 + T}$ se stejným znaménkem proto existuje $z \in \mathbb{Z}$ takové, že ${2\pi z < 1 + T < 2\pi z + \pi}$ (slovy: $1 + T$ leží v intervalu, na kterém je sinus kladný). Z Tvrzení 2 jsou zde nutně ostré nerovnosti a z Tvrzení 1 navíc platí $z \neq 0$ (z pohledu vnitřní funkce $\sin$ jsme v jiné periodě). Po odečtení $2\pi z$ obdržíme $0 < 1 + T - 2\pi z < \pi$. Označme tuto fázi $\varphi_1 \coloneqq 1 + T - 2\pi z$.

    Nyní může nastat jeden ze dvou případů: \begin{inparaenum}
        \item $\varphi_1 = 1$,
        nebo
        \item $\varphi_1 \neq 1$
    \end{inparaenum}. Oba případy vyvrátíme.

    1. Z rovnosti $\varphi_1 = 1$ dostáváme spor:
    ${\varphi_1 = 1
                \implies 1 + T - 2\pi z = 1
                \implies \pi = \frac{T + 1}{2z} \in \mathbb{Q}}$

    2. Označme $\Delta \varphi \coloneqq \varphi_1 - 1$. Protože $\varphi_1 \in (0, \pi)$ a zároveň ${\varphi_1 \neq 1}$, tak ${\Delta \varphi \in (-1, 0) \cup (0, \pi - 1)}$. Uvažujme pro teď případ $0 < \Delta \varphi < \pi - 1$. Ukážeme, že po maximálně $\ceil{\frac{\pi}{\Delta \varphi}}$ periodách se změní posloupnosti $(\sin n)^{\infty}_{n=1}$ znaménko, tedy vyvrátíme existenci periody T posloupnosti $(\sgn \sin n)^{\infty}_{n=1}$.

    Pár prvních prvků vybrané posloupnosti $(\sgn \sin (1 + n T))^{\infty}_{n = 1}$ mohou mít shodné znaménko, nicméně každou periodou se fáze ve vnitřní funkci sinus posouvá o $\Delta \varphi$. V konečném počtu $m$ `period' nastane $1 + m \cdot \Delta \varphi > \pi$ a znaménko skočí do záporných hodnot. Protože $\Delta \varphi < \pi$, nehrozí přeskočení intervalu délky $\pi$ z kladných hodnot sinu opět do kladných.

    V případě $\Delta \varphi \in (-1, 0)$ by byl důkaz analogický. Fáze by se zmenšovala až do doby, kdy po $m$ `periodách' dosáhne $1 + m \cdot \Delta \varphi < 0$ a změní se znaménko.

    \textbf{Závěr:} S předpokladem, že posloupnost $(\sgn \sin n)^{\infty}_{n=1}$ je periodická, jsme došli ke sporu. Posloupnost žádnou periodu nemá.

    \begin{figure}[H]
        \centering
        \makebox[\linewidth]{\includegraphics[width=16cm,height=3cm]{4}}
    \end{figure}
\end{proof}

\end{document}
